% No new C/D theorem is claimed for records419--427.
\subsection{An annotation on the thirteen-coin convention}

This is a check of an expository puzzle (1977-424.2), not a proposed new
solution to a class C or D research question.  The classical subject appears
in \citet{dyson1946pennies}; the publisher record was checked, but the proof
below is a self-contained reconstruction.  The difference between identifying
a counterfeit and determining its sign is essential to the notebook's number.

\begin{proposition}[Identity only, with equal pan counts]\ours
Suppose exactly one of \(m\) coins has a different mass, either heavier or
lighter, and all other coins have a common mass.  No genuine coin is supplied
externally.  Every weighing compares equal numbers of coins on an ideal
three-outcome balance.  For \(N\geq1\), the largest \(m\) whose counterfeit
identity can always be determined in at most \(N\) weighings is
\((3^N-1)/2\).  Its sign need not always be determined.
\end{proposition}

\begin{proof}
For the upper bound let the first weighing compare \(a\) coins against
\(a\), leaving \(b=m-2a\) aside.  Either unbalanced branch contains
\(2a\) possible signed counterfeits.  They must be distinguished by
\(N-1\) subsequent ternary observations, so
\(2a\leq3^{N-1}\).  Since the left side is even,
\(2a\leq3^{N-1}-1\).  On the balanced branch the \(b\) unweighed coins
have unknown signs.  For any fixed coin, its two signs can follow the same
history only if that coin is never weighed on that history; until its first
weighing both signs produce balance, whereas its first appearance produces
opposite imbalances.  Consequently only the all-balance history can fail to
determine the sign, and it can represent at most one identity.  At least
\(2b-1\) terminal histories are therefore necessary when \(b>0\).
Hence \(b\leq(3^{N-1}+1)/2\); the same bound is trivial if \(b=0\).
Adding gives \(m\leq(3^N-1)/2\).  The argument permits adaptive choices
of the subsequent weighings.

For attainment, use the alphabet \(\{-1,0,1\}\), in cyclic order
\(-1\to0\to1\to-1\).  For every nonconstant word of length \(N\),
find its first unequal adjacent pair.  Retain the word precisely when that
pair follows the cyclic order.  Negation reverses the order, so exactly one
of each opposite pair is retained.  There are
\(M=(3^N-3)/2\) retained words.  Cyclically permuting all three alphabet
symbols preserves the selection rule, so each coordinate takes each symbol
equally often among the retained words.  Their coordinate sums are therefore
zero.  Add the all-zero word, giving \(M+1=(3^N-1)/2\) columns.

Assign one column to each coin.  In weighing \(j\), put entries \(-1\)
on the left, \(1\) on the right, and leave entries \(0\) aside.  Every
weighing has equal pan counts.  A heavy counterfeit produces its column as
the outcome word, with right-heavy denoted \(1\); a light one produces the
negative column.  Distinct nonzero columns are never opposites, so these
histories identify their coins and signs.  The all-zero history identifies
the last coin, with either sign.  For \(N=1\) the retained set is empty and
the single coin is identified without weighing, as required.
\end{proof}

\begin{remark}[If the sign is required]\ours
When both identity and sign must be recovered, the balanced first branch
must distinguish all \(2b\) states, hence
\(b\leq(3^{N-1}-1)/2\).  Together with the same unbalanced-branch bound
this gives \(m\leq(3^N-3)/2\).  The nonzero retained columns above attain
that bound for \(N\geq2\).  Thus three weighings permit 12 coins with
both outputs, or 13 with identity alone.  This distinction, not a numerical
misprint, accounts for the notebook's 13.
\end{remark}

\begin{example}[Exhaustive finite checks]\oursComputed
The accompanying \texttt{coin\_check.py} generates the columns and computes
actual left and right pan masses, independently of the signed-column
formula.  For \(N=1,\ldots,6\), it checks respectively
\(1,4,13,40,121,364\) coins.  Every coordinate balances, and every outcome
determines a unique coin.  The only two-state collision is the sign of the
zero-column coin.  In the three-weighing case all 26 signed hypotheses were
enumerated, yielding 25 histories and no collision between distinct
identities.  The proof establishes the general statement; finite enumeration
only checks the implementation and the notebook's small case.
\end{example}
